\section{Teil 15}
\begin{block}[Die Trübsahlzeit]
        Mit der 70. Jahrwoche aus \bibleverse{Dan}(9:24-27) beginnt die Trübsahlzeit hier auf der Erde. Letzten Sonntag haben ich euch gesagt, dass ich glaube, dass die Entrückung vor dieser Zeit stattfinden wird. Das heisst also, wenn wir jetzt über die Trübsahlzeit sprechen sind wir als Gemeinde nicht mehr da. Es sind also nur noch Menschen auf dieser Welt die Jesus Christus als ihren Retter nicht angenommen haben. Zu Beginn dieser letzten sieben Jahre, werden die Opfer im Tempel wieder eingeführt. Aber nicht Gott ist es, der diese Opfer wieder einführt, sondern der Antichrist. Zu Beginn bringt dieser Herrscher Ruhe und Frieden auf diese Erde. Liebi nennt die ersten 3,5 Jahre die Stunde der Versuchung. 
        
        Die Menschen auf der Erde werden den Herrscher verehren und vergöttern. Er wird eine Einheitsreligion einführen. Wenn einer Auftritt und einen Frieden, oder besser gesagt eine Ruhe, auf Erden bringt, dann glaubt mir, viele  Christen die nicht Bibeltreu sind, werden auf diesen Menschen reinfallen. Diese Einheitsreligion wird nur so von Nächstenliebe strotzen, solange man der gleichen Meinung ist natürlich. Das sieht man ja schon heute, alles ist heute eine Krankheit. Ein Serienmörder hatte eine schlechte Kindheit, ein straffälliger Flüchling ist nur traumatisiert... Aber Jesus sagt in \bibleverse{Joh}(15:5), dass wir ohne ihn nichts tun können. 
        \begin{bibelbox}{SCHL}{Joh}{15:5}
            Ich bin der Weinstock, ihr seid die Reben. Wer in mir bleibt und ich in ihm, der bringt viel Frucht; denn getrennt von mir könnt ihr nichts tun.
        \end{bibelbox}
        Ein paar Menschen werden begreifen, dass etwas nicht stimmt. Wir dürfen nicht vergessen, dass wenn die Gemeinde weg ist, die Bibeln trotzdem noch da sind. Es wird also Menschen geben die sich zu Jesus hinwenden werden. In Kapitel 7 lesen wir von den 144 000 die versiegelt werden, und von einer Schar in weissen Kleidern, die aus der grossen Trübsal kommen. 
        \begin{bibelbox}{SCHL}{Offb}{7:4}
            Und ich hörte die Zahl der Versiegelten: 144000 Versiegelte, aus allen Stämmen der Kinder Israel.
        \end{bibelbox}
        und dann in den Versen 13-14;
        \begin{bibelbox}{SCHL}{Offb}{7:13-14}
            Und einer der Ältesten ergriff das Wort und sprach zu Johannes: Wer sind diese, die im weissen Kleidern bekleidet sind, und woher sind sie gekommen?

            Und ich sprach zu ihm: du weisst es! Und er sprach zu mir: Das sind die, welche aus der grossen Drangsal kommen und sie haben ihre Kleider gewaschen, und sie haben ihre Kleider weiss gemacht in dem BLut des Lammes.
        \end{bibelbox}
        In dieser Zeit werden die Schafe von den Böcken getrennt oder das Spreu vom Weizen. Wo Jesus im Gleichnis noch gesagt hat, wartet noch mit dem Unkraut schneiden, bis die Erntezeit da ist, heisst es jetzt, dass die Erntezeit da ist. Jetzt werden die Schnitter ausgesandt. \bibleverse{Mat}(13:24-30)
         \begin{bibelbox}{SCHL}{Mat}{13:30}
            Lasst beides miteinander wachsen bis zur Ernte, und zur Zeit der Ernte will ich den Schnittern sagen: Lest zuerst das Unkraut zusammen und bindet es in Bündel, dass man es verbrenne; den Weizen aber sammelt in meine Scheune.
        \end{bibelbox}
         \begin{bibelbox}{SCHL}{Offb}{14:14}
            Und ich sah, und siehe, eine weisse Wolke, und auf der Wolke sass einer, der glich einem Sohn des Menschen; er hatte auf dem Haupt eine goldene Krone und in seiner Hand eine scharfe Sichel.
        \end{bibelbox}
        In der ersten Hälfte treten die 2 Zeugen auf, die nach 3 1/2 Jahren getötet werden. Diese werden am 3. Tag auferstehen und in den Himmel aufgenommen werden. In der zweiten Hälfte der Trübsalzeit wird der Antichrist sich selbst als Gott ausgeben und sich in den Tempel setzen. Er wird die Menschen zwingen, das Malzeichen des Tieres anzunehmen, damit sie kaufen und verkaufen können. Es wird eine grosse Verfolgung geben, mit schrecklichen Gerichten. Was mich zutiefst erschreckt, ist die Tatsache, dass trotz dieser schrecklichen Ereignisse, die Menschen Gott immer noch ablehnen.
        Wenn sie nicht mal dann Gott annehmen, wenn ihnen buchstäblich der Himmel auf den Kopf fällt, wie denn jetzt in der heutigen Zeit? Naturkatastrophen? Unmöglich das die von Gott kommen, ist der Klimawandel. CO2 Steuern helfen. Wir aber sollen und müssen den Menschen von Jesus erzählen, aber wir können sie nicht zwingen. Am besten ist es, mit unserem Leben Zeugnis geben. 
        
    \end{block}